% ID: SEL-D0025-OA-E02
% ESTADO: cotejada. Segunda pasada visual completa y borrador compilado revisados.
% Fuente: Selectividad/matematicas_aplicadas/suplemente_junio_EXAMEN 4-ANDALUCÍA.pdf
% SHA-256: c1a3a7a37d38b04359115c64762b1f8a543e4d66155338931de84ef6c304da2c
% Página: 1; opción A; ejercicio original 2.

En una especie animal la contracción del iris, en décimas de milímetro, después de exponer el ojo a una luz brillante durante un determinado tiempo, viene dada por
\[
f(t)=\begin{cases}
t^2 & \text{si }0\leq t\leq 2\\[4pt]
\dfrac{4}{t-1} & \text{si }t>2
\end{cases}
\]
donde $t$ es el tiempo, en segundos, que transcurre desde que se concentra la luz en el ojo.
\begin{enumerate}[label=\alph*)]
\item \textbf{(1 punto)} Estudie la continuidad y la derivabilidad de la función $f$.
\item \textbf{(1 punto)} Represente gráficamente la función $f$, determinando los intervalos de crecimiento y decrecimiento y sus asíntotas, en caso de que existan.
\item \textbf{(0.5 puntos)} Determine en qué instante se obtiene la máxima contracción y su valor.
\end{enumerate}
